\section{Proofs for the Weak-Convergence Exercises}
\label{sec:annex-p15-weak-convergence-exercises}

\begin{annexitem}{annex:p15.wce-dyadic-densities}{thm:p15-wce-dyadic-densities}{Exact cell integrals and pointwise failure}
\begin{proof}
Each of the \(2^n\) half-cells in \(E_n\) has length
\(2^{-n-1}\), so \(f_n\geq0\) and
\(\int f_n=2\cdot2^n\cdot2^{-n-1}=1\).
Put \(\ell=2^{-n}\). At every cell endpoint \(j\ell\),
\(F_n(j\ell)=j\ell=F(j\ell)\).
For \(x=j\ell+r\), \(0\leq r<\ell\),
\[
 F_n(x)-F(x)=
 \begin{cases}
  r,&0\leq r\leq\ell/2,\\
  \ell-r,&\ell/2\leq r<\ell.
 \end{cases}
 \]
Outside \([0,1]\) both CDFs are zero or one.
Thus their uniform difference is exactly \(\ell/2\), proving
the weak limit.
At every \(x\in[0,1]\), including \(x=1\) where \(f_n(1)=0\),
the density value is either zero or two. Its distance from one
is exactly one.
The integral of the absolute density difference is therefore one.
For any measurable set \(A\), the difference between the two
probabilities is the integral over \(A\) of
\(f_n-\mathbf1_{[0,1]}\). Its positive and negative parts
each have integral \(1/2\), so the absolute difference is at
most \(1/2\); choosing \(A=E_n\) attains that bound.
Thus the total-variation distance is exactly \(1/2\).
This is consistent with weak convergence,
which does not require either of those stronger forms of convergence.
\end{proof}
\end{annexitem}

\begin{annexitem}{annex:p15.wce-extreme-maxima}{thm:p15-wce-extreme-maxima}{The three finite CDF formulas and their limits}
\begin{proof}
The event \(M_n\leq x\) is the intersection of the independent
events \(X_j\leq x\). Thus its probability is \(F(x)^n\).
For the Pareto case and fixed \(y>0\), the threshold
\(n^{1/\alpha}y\) is at least one for all sufficiently large \(n\);
at those indices
\[
 \mathbb P(M_n/n^{1/\alpha}\leq y)
                         =(1-y^{-\alpha}/n)^n.
 \]
This tends to \(e^{-y^{-\alpha}}\) by the proved elementary
logarithm limit. For \(y\leq0\) the probability is zero.
The limiting expression increases continuously from zero at
\(y=0\) to one at infinity.

For the reversed Weibull case and fixed \(y<0\), the threshold
\(y/n^{1/\beta}\) lies in \([-1,0)\) eventually. Therefore
\[
 \mathbb P(n^{1/\beta}M_n\leq y)=(1-|y|^\beta/n)^n
                                  \longrightarrow e^{-|y|^\beta}.
 \]
For \(y\geq0\) its probability is one. The resulting CDF
is continuous at zero, tends to zero at negative infinity,
and is nondecreasing.

For the exponential case and any fixed real \(y\),
\(\log n+y\geq0\) eventually, giving
\[
 \mathbb P(M_n-\log n\leq y)=(1-e^{-y}/n)^n
                                  \longrightarrow e^{-e^{-y}}.
 \]
This last continuous function has limits zero and one at the
two infinities. Each limiting function is therefore a proper
CDF, and convergence at every point proves weak convergence.
Every finite formula was used only inside its actual support.

All three finite maximum CDFs are continuous and strictly
increasing in the interiors of their supports.
Solving \(F(x)^n=u\), \(0<u<1\), gives exactly the three
quantiles stated. Since
\[
 n(1-u^{1/n})=n(1-e^{(\log u)/n})\longrightarrow-\log u>0,
 \]
their respective normalized limits are
\((-\log u)^{-1/\alpha}\), \(-(-\log u)^{1/\beta}\), and
\(-\log(-\log u)\). These are precisely the corresponding
limiting CDF quantiles.
\end{proof}
\end{annexitem}

\begin{annexitem}{annex:p15.wce-normal-maxima}{thm:p15-wce-normal-maxima}{Tail ratios, quantile centering and the first-order scale}
\begin{proof}
The Gaussian tail bounds already supplied in
Theorem~\ref{thm:p15-stirling-tail-premises} imply
\[
 \overline\Phi(x)\sim\phi(x)/x\quad(x\to\infty).
 \]
They are the needed bounds of source Theorem 1.2.6.
For fixed \(\theta\), let \(z=x+\theta/x\).
Then \(z\to\infty\), \(z/x\to1\), and
\(z^2-x^2=2\theta+\theta^2/x^2\).
Consequently
\[
 \frac{\overline\Phi(z)}{\overline\Phi(x)}
 \sim\frac{x}{z}\exp\{-(z^2-x^2)/2\}
                         \longrightarrow e^{-\theta}.
 \]
This proves the first requested ratio for every real \(\theta\).

The normal density is strictly positive and continuous.
Thus \(\overline\Phi\) is continuous and strictly decreasing
from one to zero, with value \(1/2\) at zero.
For \(n\geq3\), the intermediate value theorem gives the
unique positive solution \(b_n\).
For every fixed \(B\), eventually \(1/n<\overline\Phi(B)\),
so \(b_n>B\). Hence \(b_n\to\infty\).
For fixed \(y\in\mathbb R\), the first ratio gives
\[
 n\overline\Phi(b_n+y/b_n)
  =\frac{\overline\Phi(b_n+y/b_n)}{\overline\Phi(b_n)}
                                      \longrightarrow e^{-y}.
 \]
If \(q_n=\overline\Phi(b_n+y/b_n)\), then \(q_n\to0\),
\(nq_n\to e^{-y}\), and \(nq_n^2\to0\).
The elementary bound
\(|\log(1-q)+q|\leq q^2\) for \(0\leq q\leq1/2\)
therefore gives
\[
 \mathbb P\{b_n(M_n-b_n)\leq y\}
  =(1-q_n)^n\longrightarrow e^{-e^{-y}}.
 \]

The tail asymptotic at \(b_n\), followed by its logarithm,
gives the precise relation
\[
 \log n=\frac{b_n^2}{2}+\log b_n
                         +\frac12\log(2\pi)+o(1).
 \]
Since \(b_n\to\infty\) and \(\log b_n/b_n^2\to0\), division
by \(b_n^2/2\) yields \(b_n\sim\sqrt{2\log n}\).
Then \(\log b_n=\tfrac12\log(2\log n)+o(1)\);
substitution into the same relation gives
\[
 b_n^2=2\log n-\log\log n-\log(4\pi)+o(1).
 \]
Finally \(Z_n=b_n(M_n-b_n)\) converges to a proper law,
so its laws are tight. For every \(\varepsilon>0\),
\(\mathbb P(|Z_n|>\varepsilon b_n^2)\to0\), by first
bounding the tails at a fixed large threshold and then using
\(b_n^2\to\infty\). Thus
\[
 M_n/b_n=1+Z_n/b_n^2\longrightarrow1
                                      \quad\text{in probability}.
 \]
Multiplication by \(b_n/\sqrt{2\log n}\to1\) proves the final
ratio assertion. This maximum limit used normal tail estimates,
not a central limit theorem for sums.
\end{proof}
\end{annexitem}

\begin{annexitem}{annex:p15.wce-bounded-distance}{thm:p15-wce-bounded-distance}{The whole bounded metric and both comparison bounds}
\begin{proof}
Let \(h(t)=t/(1+t)\), \(t\geq0\).
It is nonnegative, increasing, bounded by one, and zero exactly
at zero. Thus \(\beta(X,Y)\) is finite, symmetric, unchanged
on almost-sure classes, and zero exactly when \(X=Y\) almost surely.
For \(s,t\geq0\),
\[
 h(s+t)=\frac{s}{1+s+t}+\frac{t}{1+s+t}
                                  \leq h(s)+h(t).
 \]
The pointwise triangle inequality and monotonicity of \(h\)
give \(h(|X-Z|)\leq h(|X-Y|)+h(|Y-Z|)\).
Taking expectations proves the metric triangle inequality.

Put \(D=|X-Y|\) and \(a=\alpha(X,Y)\).
The existing Ky Fan theorem gives \(0\leq a\leq1\) and
\(\mathbb P(D>a)\leq a\), including the attained endpoint.
For \(0\leq r<a\), the definition of the infimum instead gives
\(\mathbb P(D>r)>r\).
Hence for such \(r\),
\[
 \beta(X,Y)=\mathbb Eh(D)
             \geq h(r)\mathbb P(D>r)\geq r^2/(1+r).
 \]
Let \(r\uparrow a\). If \(a=0\), the same lower bound
is simply nonnegativity.
For the upper bound, put \(p=\mathbb P(D>a)\leq a\).
Splitting the expectation gives
\[
 \mathbb Eh(D)\leq (1-p)h(a)+p
                          \leq (1-a)h(a)+a.
 \]
This is the exact source inequality.
Because \(a^2/(1+a)\geq a^2/2\) and
\(2a/(1+a)\leq2a\), the inequalities show
\(\beta(X_n,X)\to0\) if and only if \(\alpha(X_n,X)\to0\).
The proved Ky Fan characterization makes this equivalent to
convergence in probability. This supplies the entire referred
Exercise 2.3.6 premise.

For the lower-bound example, take \(X=0\) and let \(Y=a\)
with probability \(a\), zero otherwise.
For \(0\leq r<a\), its tail above \(r\) is \(a>r\);
at \(r=a\) it is zero. Thus its Ky Fan distance is \(a\)
and its bounded distance is \(a^2/(1+a)\).
For the upper bound, let \(Y=M>a\) with probability \(a\)
and \(Y=a\) otherwise. For \(r<a\) its tail is one,
and at \(r=a\) its tail is \(a\).
The Ky Fan distance is again \(a\), while
\(\beta=(1-a)h(a)+a h(M)\) tends to the upper bound as
\(M\to\infty\). This also handles \(a=1\).
At \(a=0\), take the identical-zero variables.
Every example uses finite random variables.
\end{proof}
\end{annexitem}

\begin{annexitem}{annex:p15.wce-integer-laws}{thm:p15-wce-integer-laws}{Mass convergence and the complete tail argument}
\begin{proof}
Write \(p_n(m)=\mathbb P(X_n=m)\), and use subscript
\(\infty\) for the limiting law. At each noninteger point
the CDF of the integer-valued limit is continuous. Thus weak
convergence gives
\[
 p_n(m)=F_n(m+\tfrac12)-F_n(m-\tfrac12)
        \longrightarrow
        F_\infty(m+\tfrac12)-F_\infty(m-\tfrac12)=p_\infty(m).
 \]

Conversely suppose the masses converge at every integer.
For arbitrary \(\varepsilon>0\), choose a finite set
\(A\subseteq\mathbb Z\) with
\(\sum_{m\notin A}p_\infty(m)<\varepsilon\).
This is possible because the limiting masses sum to one.
The sum of errors on \(A\) tends to zero, and
\(\sum_{m\notin A}p_n(m)=1-\sum_{m\in A}p_n(m)\)
converges to the corresponding limiting tail.
Therefore
\[
 \limsup_n\sum_{m\in\mathbb Z}|p_n(m)-p_\infty(m)|
 \leq2\varepsilon.
 \]
Let \(\varepsilon\downarrow0\).
The positive and negative parts of a summable mass difference
with total sum zero have equal sums. Taking the event of all
indices where the difference is nonnegative shows
\[
 \|\mathcal L(X_n)-\mathcal L(X_\infty)\|_{\mathrm{TV}}
       =\frac12\sum_m|p_n(m)-p_\infty(m)|\longrightarrow0.
 \]
In particular every CDF difference is bounded by this
total-variation distance, proving weak convergence.
\end{proof}
\end{annexitem}

\begin{annexitem}{annex:p15.wce-spherical-area}{thm:p15-wce-spherical-area}{Volume, area and the polar identity from cones}
\begin{proof}
Product Lebesgue measure gives
\[
 (2/\sqrt d)^d
   =\lambda_d([-1/\sqrt d,1/\sqrt d]^d)
   \leq v_d\leq\lambda_d([-1,1]^d)=2^d.
 \]
Thus \(v_d\) is finite and positive. The origin has measure zero.
The radial projection \(\pi(x)=x/\|x\|\) is continuous on
\(\mathbb R^d\setminus\{0\}\).
Consequently \(C(A)=\pi^{-1}(A)\cap\{0<\|x\|\leq1\}\)
is Borel. Cones of disjoint angular sets are disjoint, so
countable additivity of volume proves countable additivity
of \(\sigma_d\), and \(\sigma_d(\mathbb S^{d-1})=1\).

We prove the volume invariances used here.
For a translation \(b\in\mathbb R^d\), the measures
\(A\mapsto\lambda_d(A+b)\) and \(\lambda_d(A)\) agree on
half-open axis rectangles. For a positive scalar \(t\),
\(A\mapsto t^{-d}\lambda_d(tA)\) and \(\lambda_d(A)\)
also agree on those rectangles. All these measures are
sigma-finite and the rectangles generate the Borel sigma-algebra.
Measure uniqueness therefore gives translation invariance and
\(\lambda_d(tA)=t^d\lambda_d(A)\) for every Borel \(A\).

For orthogonal \(O\), put \(\nu(A)=\lambda_d(O^{-1}A)\).
This measure is translation invariant and finite on bounded sets,
since \(O\) preserves Euclidean norm.
Any such measure is a constant multiple of \(\lambda_d\):
if \(c=\nu([0,1)^d)\), partitioning the unit cube into
\(m^d\) translated half-open subcubes gives each mass \(c/m^d\).
Tiling rectangles with side lengths \(k_j/m\) gives mass
\(c\prod_j(k_j/m)\). Approximate arbitrary positive side lengths
from below by rational ones and use continuity from below;
translation then gives the same formula for all half-open
axis rectangles. Sigma-finite measure uniqueness gives
\(\nu=c\lambda_d\) on every Borel set.
But \(O^{-1}B_d=B_d\), so
\(v_d=\nu(B_d)=cv_d\), forcing \(c=1\).
This proves orthogonal invariance of volume.
Since \(O C(A)=C(OA)\), it also proves orthogonal invariance
of the cone probability.

For \(0\leq a<b<\infty\), the radial slab of angular set
\(A\) has the exact volume
\[
 \lambda_d\{x:a<\|x\|\leq b,\ \pi(x)\in A\}
   =(b^d-a^d)\lambda_d(C(A))
   =d v_d\,\sigma_d(A)\int_a^b r^{d-1}\,dr.
 \]
It is the difference of two nested dilated cones;
their common missing origin has zero volume.
Taking \(a=r,b=r+\varepsilon\), dividing by
\(\varepsilon\), and letting \(\varepsilon\downarrow0\)
gives the claimed surface-area measure on every Borel patch.
In particular it is a countably additive measure and its
normalization is the stated cone probability at every radius.
For \(d=1\), each singleton cone is an interval of length one,
giving area one at each of the two sphere points.

For the polar integration formula, push \(\lambda_d\) on the
punctured space forward by \(x\mapsto(\|x\|,\pi(x))\).
The preceding slab identity identifies this sigma-finite measure,
on every product rectangle \((a,b]\times A\), with the product of
\(d v_d r^{d-1}\,dr\) and \(\sigma_d\).
These rectangles generate the product Borel sigma-algebra.
Indeed both factors have countable topological bases, and
the inverse map \((r,u)\mapsto ru\) is continuous.
The sigma-finite uniqueness theorem therefore identifies the
measures on all Borel sets. Integration of nonnegative simple
functions and monotone approximation gives the asserted formula
for every nonnegative Borel \(h\). The positive/negative
decomposition gives its integrable signed version.

Finally Tonelli and the one-dimensional Gaussian integral
of Theorem~\ref{thm:p15-normal-normalization} give
\(\int_{\mathbb R^d}e^{-\|x\|^2/2}\,dx=(2\pi)^{d/2}\).
The polar formula and the substitution \(t=r^2/2\) give
\[
 (2\pi)^{d/2}
 =d v_d\int_0^\infty e^{-r^2/2}r^{d-1}\,dr
 =d v_d\,2^{d/2-1}\Gamma(d/2).
 \]
Thus \(d v_d=2\pi^{d/2}/\Gamma(d/2)\).
The previously proved gamma recurrence
\(\Gamma(d/2+1)=(d/2)\Gamma(d/2)\) gives the formula for \(v_d\).
This derivation used no multivariable substitution theorem and
did not presume a Gaussian definition of uniform surface measure.
\end{proof}
\end{annexitem}

\begin{annexitem}{annex:p15.wce-sphere-coordinate}{thm:p15-wce-sphere-coordinate}{The normalized Gaussian vector has the prescribed area law}
\begin{proof}
For a fixed \(n\), let \(Y=(Y_1,\ldots,Y_n)\),
\(R=\|Y\|\), and \(U=Y/R\) on \(R>0\).
The origin has probability zero because the product normal law
has a Lebesgue density. Define \(U=(1,0,\ldots,0)\) there.
The joint density of \(Y\) is
\[
 (2\pi)^{-n/2}e^{-\|y\|^2/2}.
 \]
For a Borel radial set \(B\subseteq(0,\infty)\) and
angular set \(A\subseteq\mathbb S^{n-1}\), polar integration gives
\[
 \mathbb P(R\in B,U\in A)
 =n v_n(2\pi)^{-n/2}\sigma_n(A)
                    \int_B e^{-r^2/2}r^{n-1}\,dr.
 \]
With \(B=(0,\infty)\), the factor multiplying \(\sigma_n(A)\)
is one, since the normal density is normalized.
Thus \(U\) has exactly the geometric surface probability
\(\sigma_n\), and the same rectangle identity proves that
\(R\) and \(U\) are independent.
Multiplication by \(\sqrt n\) gives the uniform surface
probability on the sphere of radius \(\sqrt n\), including
the \(n=1\) two-point case.

On the common infinite Gaussian product space, the integrable
strong law (Theorem~\ref{thm:p15-etemadi-strong-law}) applied
to \(Y_j^2\), whose mean is one, gives
\[
 D_n=\frac1n\sum_{j=1}^nY_j^2\longrightarrow1
                                            \quad\text{almost surely}.
 \]
Consequently \(Y_1/\sqrt{D_n}\to Y_1\) almost surely
where the displayed fraction is eventually defined.
The exceptional finite-index cases were assigned sphere values
as in the statement and have probability zero.
Almost-sure convergence implies convergence in probability
(\cref{thm:p15-w02b-p15-013-as-probability}~(1))
and hence weak convergence, by the probability bridge of
\cref{thm:p15-weak-metric-premises}.
Since each coupled vector has the required surface law, this
proves the exercise for any sequence with those marginal laws.

For its moments, the surface law is unchanged by coordinate
reflections and permutations. Boundedness of every coordinate
therefore gives \(\mathbb EX^{(n)}_1=0\); summing the
squared coordinates, which total \(n\), gives
\(\mathbb E(X^{(n)}_1)^2=1\).
For the fourth moment, one-dimensional integration by parts
using \(\phi'(t)=-t\phi(t)\) gives
\[
 \mathbb EY_1^4=-\int_{\mathbb R}t^3\phi'(t)\,dt
       =3\int_{\mathbb R}t^2\phi(t)\,dt=3.
 \]
The boundary term \(t^3\phi(t)\) vanishes at both infinities;
all integrals are finite by Gaussian decay.
Independence yields
\[
 \mathbb ER^4
 =\mathbb E\left(\sum_{j=1}^nY_j^2\right)^2
 =3n+n(n-1)=n(n+2).
 \]
The already proved independence of \(R,U\) and \(Y_1=R U_1\)
then gives
\(\mathbb EU_1^4=3/[n(n+2)]\).
Multiplication by the fourth power of \(\sqrt n\) proves
\(\mathbb E(X^{(n)}_1)^4=3n/(n+2)\).
\end{proof}
\end{annexitem}

\begin{annexitem}{annex:p15.wce-moment-concentration}{thm:p15-wce-moment-concentration}{The strict convexity gap}
\begin{proof}
Put \(p=\beta/\alpha>1\) and \(Z_n=Y_n^\alpha\).
The function \(D(z)=z^p-pz+p-1\), \(z\geq0\), has
\(D(1)=0\) and derivative \(p(z^{p-1}-1)\).
It decreases on \([0,1]\), increases on \([1,\infty)\),
and is strictly positive except at one.
For \(0<\varepsilon<1\), the monotonicity of \(y^\alpha\)
therefore gives
\[
 D(Y_n^\alpha)\geq c_\varepsilon
                    \mathbf1_{\{|Y_n-1|>\varepsilon\}},
 \quad
 c_\varepsilon=
 \min\{D((1-\varepsilon)^\alpha),D((1+\varepsilon)^\alpha)\}>0.
 \]
Convergence of both nonnegative moments to finite values makes
them finite for all sufficiently large \(n\).
At those indices expectation and finite linearity give
\[
 \mathbb ED(Y_n^\alpha)
   =\mathbb EY_n^\beta-p\mathbb EY_n^\alpha+p-1\to0.
 \]
The displayed lower bound yields the quantitative estimate and
convergence in probability for every \(0<\varepsilon<1\).
Any larger deviation event is contained in one of these smaller
deviation events, so all positive tolerances are covered.
Nothing in the proof uses independence or a first moment of \(Y_n\).
\end{proof}
\end{annexitem}

\begin{annexitem}{annex:p15.wce-fourth-moment}{thm:p15-wce-fourth-moment}{The complete fourth-moment lower bound}
\begin{proof}
Cauchy--Schwarz gives
\((\mathbb EX^2)^2\leq\mathbb EX^4\), so the assumptions
are impossible when \(K<1\).
If \(K\geq1\) and \(0\leq y<1\), split at the strict tail:
\[
 \begin{split}
 1=\mathbb EX^2
 &=\mathbb E(X^2;|X|\leq y)+\mathbb E(X^2;|X|>y)\\
 &\leq y^2+
         (\mathbb EX^4)^{1/2}\mathbb P(|X|>y)^{1/2}\\
 &\leq y^2+\sqrt K\,\mathbb P(|X|>y)^{1/2}.
 \end{split}
 \]
Subtract \(y^2<1\) and square the nonnegative inequality
to obtain the bound. For \(y<0\), nonnegativity of \(|X|\)
makes the event certain.
At \(y=0\), take a variable with probability \(1/(2K)\)
at each of \(\sqrt K,-\sqrt K\) and probability \(1-1/K\)
at zero. Its second moment is one, its fourth moment is \(K\),
and \(\mathbb P(|X|>0)=1/K\), proving attainment.
\end{proof}
\end{annexitem}

\begin{annexitem}{annex:p15.wce-boundaries}{thm:p15-wce-boundaries}{Concrete failures when the assumptions are removed}
\begin{proof}
For a single finite Pareto variable \(X\), take \(X_j=X\)
at every index. Each marginal is correct but \(M_n=X\),
so \(M_n/n^{1/\alpha}\to0\) almost surely.
The independent limiting CDF is not the point mass at zero.

The deterministic real variables \(1/n\) converge to zero
pointwise and hence weakly, but their probability at zero
is always zero instead of the limiting mass one.
For integer constants \(X_n=n\), each fixed integer mass
eventually vanishes. The candidate zero mass function sums
to zero. There is no proper weak limit, since these laws
are not tight: for every fixed \(M\), their mass outside
\([-M,M]\) is one eventually.

The two-valued \(Y\) in the statement has
\(\mathbb EY^\alpha=(0+2)/2=1\).
Its probability of being zero is always \(1/2\), so it
cannot converge in probability to one.
With \(Y_n=-1\), the second and fourth moments are both one,
whereas \(\mathbb P(|Y_n-1|>1)=1\).
This is a concrete failure of the nonnegativity hypothesis
with powers that are well-defined at the negative value.

For the last family,
\(\mathbb EX_n^2=n(1/n)=1\) and
\(\mathbb EX_n^4=n^2(1/n)=n\).
For every fixed \(0\leq y<1\), eventually
\(\mathbb P(|X_n|>y)=1/n\to0\).
Thus a bounded second moment alone cannot supply the source's
uniform positive lower bound.
The fair-sign variable has \(|X|=1\), giving zero probability
of the strict event at threshold one despite both specified
moments equalling one.
\end{proof}
\end{annexitem}
